\newcounter{stoppagebreaks}
\documentclass{article}

% packages
\usepackage{amsmath, amsthm, thmtools, amsfonts, amssymb, luacode, catchfile, hyperref, ifthen, graphicx, bbm}
\ifcsname c@kobocompile\endcsname
	\usepackage[a5paper, total={1072pt, 1448pt}, margin=10pt, includeheadfoot]{geometry} % set page margins
\else
	\usepackage[a4paper, margin=50pt, includeheadfoot]{geometry}
\fi
\usepackage[shortlabels]{enumitem}
\usepackage[skip=3pt, indent=0pt]{parskip}

% language
\usepackage[bidi=basic, layout=tabular, provide=*]{babel}
\ifcsname c@english\endcsname
	\babelprovide[main, import]{english}
\else
	\babelprovide[main, import, alph=letters.gershayim]{hebrew}
	\setlist[enumerate,2]{label={\alph*.}}
	\babelprovide{rl}
\fi
\babelfont{rm}[Renderer=Harfbuzz]{Libertinus Serif}
\babelfont{sf}{Libertinus Sans}
\babelfont{tt}{Libertinus Mono}

% style
\ifcsname c@stoppagebreaks\endcsname
\else
	\AddToHook{cmd/section/before}{\clearpage}	% Add line break before section
\fi
\linespread{1.3}
\setcounter{secnumdepth}{0}		% Remove default number tags from sections, this won't do well with theorems
\AtBeginDocument{\setlength{\belowdisplayskip}{3pt}}
\AtBeginDocument{\setlength{\abovedisplayskip}{3pt}}
\graphicspath{ {../images/} }

% operators
\DeclareMathOperator\cis{cis}
\DeclareMathOperator\Sp{Span}
\DeclareMathOperator\tr{tr}
\DeclareMathOperator\im{Im}
\DeclareMathOperator\re{Re}
\DeclareMathOperator\diag{diag}
\DeclareMathOperator*\lowlim{\underline{lim}}
\DeclareMathOperator*\uplim{\overline{lim}}
\DeclareMathOperator\rng{rng}
\DeclareMathOperator\Sym{Sym}
\DeclareMathOperator\Arg{Arg}
\DeclareMathOperator\Log{Log}
\DeclareMathOperator\dom{dom}
\DeclareMathOperator\supp{Supp}
\DeclareMathOperator\var{Var}
\DeclareMathOperator\cov{Cov}
\DeclareMathOperator\End{End}
\DeclareMathOperator\id{id}
\DeclareMathOperator\cl{cl}
\DeclareMathOperator\gal{Gal}
\DeclareMathOperator\dist{dist}
\DeclareMathOperator\form{form}
\DeclareMathOperator\term{term}
\DeclareMathOperator\sent{sent}
\DeclareMathOperator\Th{Th}
\DeclareMathOperator\Hom{Hom}
\DeclareMathOperator\Aut{Aut}
\DeclareMathOperator\Char{char}
\DeclareMathOperator\vol{vol}
\DeclareMathOperator\GL{GL}
\DeclareMathOperator\SL{SL}
\DeclareMathOperator\lcm{lcm}

% commands
\newcommand{\cupdot}{\mathbin{\mathaccent\cdot\cup}}
\newcommand{\indicator}{\mathbbm{1}}
\newcommand{\AAA}[0]{\mathbb{A}}
\newcommand{\Aa}[0]{\mathcal{A}}
\newcommand{\BB}[0]{\mathbb{B}}
\newcommand{\Bb}[0]{\mathcal{B}}
\newcommand{\CC}[0]{\mathbb{C}}
\newcommand{\Cc}[0]{\mathcal{C}}
\newcommand{\DD}[0]{\mathbb{D}}
\newcommand{\Dd}[0]{\mathcal{D}}
\newcommand{\EE}[0]{\mathbb{E}}
\newcommand{\FF}[0]{\mathbb{F}}
\newcommand{\Ff}[0]{\mathcal{F}}
\newcommand{\Ii}[0]{\mathcal{I}}
\newcommand{\Gg}[0]{\mathcal{G}}
\newcommand{\HH}[0]{\mathbb{H}}
\newcommand{\Hh}[0]{\mathcal{H}}
\newcommand{\Kk}[0]{\mathcal{K}}
\newcommand{\Ll}[0]{\mathcal{L}}
\newcommand{\Mm}[0]{\mathcal{M}}
\newcommand{\NN}[0]{\mathbb{N}}
\newcommand{\Nn}[0]{\mathcal{N}}
\newcommand{\PP}[0]{\mathbb{P}}
\newcommand{\Pp}[0]{\mathcal{P}}
\newcommand{\QQ}[0]{\mathbb{Q}}
\newcommand{\RR}[0]{\mathbb{R}}
\newcommand{\Rr}[0]{\mathcal{R}}
\newcommand{\Ss}[0]{\mathcal{S}}
\newcommand{\TT}[0]{\mathbb{T}}
\newcommand{\Tt}[0]{\mathcal{T}}
\newcommand{\UU}[0]{\mathbb{U}}
\newcommand{\Uu}[0]{\mathcal{U}}
\newcommand{\Vv}[0]{\mathcal{V}}
\newcommand{\Ww}[0]{\mathcal{W}}
\newcommand{\Xx}[0]{\mathcal{X}}
\newcommand{\Yy}[0]{\mathcal{Y}}
\newcommand{\ZZ}[0]{\mathbb{Z}}
\newcommand{\Zz}[0]{\mathcal{Z}}
\newcommand{\acts}[0]{\circlearrowright}
\newcommand{\explain}[2] {
	\begin{flalign*}
		 && \text{#2} && \text{#1}
	\end{flalign*}
}
\newcommand{\maketitleprint}[1][green]{
	\begin{center}
		\includegraphics[width=6cm]{dragon_#1}
	\end{center}
}

% theorem commands
\newtheoremstyle{c_remark}
	{}	% Space above
	{}	% Space below
	{}% Body font
	{}	% Indent amount
	{\bfseries}	% Theorem head font
	{}	% Punctuation after theorem head
	{.5em}	% Space after theorem head
	{\thmname{#1}\thmnumber{ #2}\thmnote{ \normalfont{\text{(#3)}}}}	% head content
\newtheoremstyle{c_definition}
	{3pt}	% Space above
	{3pt}	% Space below
	{}% Body font
	{}	% Indent amount
	{\bfseries}	% Theorem head font
	{}	% Punctuation after theorem head
	{.5em}	% Space after theorem head
	{\thmname{#1}\thmnumber{ #2}\thmnote{ \normalfont{\text{(#3)}}}}	% head content
\newtheoremstyle{c_plain}
	{3pt}	% Space above
	{3pt}	% Space below
	{\itshape}% Body font
	{}	% Indent amount
	{\bfseries}	% Theorem head font
	{}	% Punctuation after theorem head
	{.5em}	% Space after theorem head
	{\thmname{#1}\thmnumber{ #2}\thmnote{ \text{(#3)}}}	% head content

\ifcsname c@english\endcsname
	\theoremstyle{plain}
	\newtheorem{theorem}{Theorem}[section]
	\newtheorem*{theorem*}{Theorem}
	\newtheorem{lemma}[theorem]{Lemma}
	\newtheorem{proposition}[theorem]{Proposition}
	\newtheorem*{proposition*}{Proposition}

	\theoremstyle{definition}
	\newtheorem{definition}[theorem]{Definition}
	\newtheorem*{definition*}{Definition}
	\newtheorem{example}{Example}[section]
	\newtheorem{exercise}{Exercise}[section]

	\theoremstyle{remark}
	\newtheorem*{remark}{Remark}
	\newtheorem*{solution}{Solution}
	\newtheorem{corollary}[theorem]{Corollary}
	\newtheorem{notation}[theorem]{Notation}
\else
	\theoremstyle{c_plain}
	\newtheorem{theorem}{משפט}[section]
	\newtheorem*{theorem*}{משפט}
	\newtheorem{lemma}[theorem]{למה}
	\newtheorem{proposition}[theorem]{טענה}
	\newtheorem*{proposition*}{טענה}

	\theoremstyle{c_definition}
	\newtheorem{definition}[theorem]{הגדרה}
	\newtheorem*{definition*}{הגדרה}
	\newtheorem{example}{דוגמה}[section]
	\newtheorem{exercise}{תרגיל}[section]

	\theoremstyle{c_remark}
	\newtheorem*{remark}{הערה}
	\newtheorem*{solution}{פתרון}
	\newtheorem{corollary}[theorem]{מסקנה}
	\newtheorem{notation}[theorem]{סימון}
\fi

% Questions related commands
\newcounter{question}
\setcounter{question}{1}
\newcounter{sub_question}
\setcounter{sub_question}{1}

\ifcsname c@english\endcsname
	\newcommand{\question}[1][0]{
		\ifthenelse{#1 = 0}{}{\setcounter{question}{#1}}
		\section{Question \arabic{question}}
		\addtocounter{question}{1}
		\setcounter{sub_question}{1}
	}

	\newcommand{\subquestion}[1][0]{
		\ifthenelse{#1 = 0}{}{\setcounter{sub_question}{#1}}
		\subsection{Part \alph{sub_question}}
		\addtocounter{sub_question}{1}
	}
\else
	\newcommand{\question}[1][0]{
		\ifthenelse{#1 = 0}{}{\setcounter{question}{#1}}
		\section{שאלה \arabic{question}}
		\addtocounter{question}{1}
		\setcounter{sub_question}{1}
	}

	\newcommand{\subquestion}[1][0]{
		\ifthenelse{#1 = 0}{}{\setcounter{sub_question}{#1}}
		\subsection{סעיף \localecounter{letters.gershayim}{sub_question}}
		\addtocounter{sub_question}{1}
	}
\fi

% import lua and start of document
\directlua{common = require ('../common')}

\GetEnv{AUTHOR}

% headers
\author{\AUTHOR}
\date\today

\setlist[enumerate,1]{label={\alph*.}}
% chktex-file 9
% chktex-file 17

\usepackage{pgfplots}
\pgfplotsset{compat=1.18}

\title{מתמטיקה א' --- פתרון תרגיל 2}
\date{}

\begin{document}

\maketitle

\section{שאלה 1}
להלן הזוגות המתאימים של הקבוצות השוות (בקבוצות, סדר האיברים וכפילויות אינם משנים):
\begin{itemize}
    \item $\{\} = \emptyset$
    \item $\{\text{סוס}, \text{כלב}, \text{דג}\} = \{\text{דג}, \text{כלב}, \text{סוס}, \text{סוס}\}$
    \item $\{a, a, b, b, c\} = \{a, b, c\}$
    \item $\{1, 3, 2\} = \{1, 2, 2, 3\}$
\end{itemize}

\section{שאלה 2}
נתונה הקבוצה $A = \{1, 2, \{1\}, \{3\}\}$:
\begin{enumerate}
    \item $1 \in A$ --- \textbf{נכון} (האיבר 1 נמצא בקבוצה).
    \item $3 \in A$ --- \textbf{לא נכון} (האיבר 3 לא בקבוצה, אלא הקבוצה $\{3\}$).
    \item $\emptyset \in A$ --- \textbf{לא נכון} (הקבוצה הריקה אינה איבר ב־$A$).
    \item $\emptyset \subseteq A$ --- \textbf{נכון} (הקבוצה הריקה היא תת־קבוצה של כל קבוצה).
    \item $\{2\} \in A$ --- \textbf{לא נכון} (הקבוצה $\{2\}$ אינה איבר ב־$A$, אלא המספר 2).
    \item $\{2\} \subseteq A$ --- \textbf{נכון} (האיבר 2 נמצא ב־$A$, ולכן הקבוצה המכילה אותו היא תת־קבוצה).
    \item $\{\{1\}, \{3\}\} \subseteq A$ --- \textbf{נכון} (האיברים $\{1\}$ ו־$\{3\}$ שניהם שייכים ל־$A$).
    \item $\{1, 2, 3\} = A$ --- \textbf{לא נכון} (האיברים שונים, למשל 3 לא ב־$A$).
\end{enumerate}

\section{שאלה 3}
פתרונות לפעולות בין קבוצות וקטעים:
\begin{enumerate}
    \item $[-4, 7] \cap [4, 7] = [4, 7]$
    \item $[-4, 7] \cup [4, 7] = [-4, 7]$
    \item $(2, 5] \cap [4, 5) = [4, 5)$
    \item $\emptyset \cap (1, 2) = \emptyset$
    \item $\emptyset \cup (3, 5) = (3, 5)$
    \item $\{3, 4, 7\} \cap \{1, 2, 3\} = \{3\}$
    \item $\{1, 2, 3\} \setminus \{1, 2, 3, 4\} = \emptyset$
\end{enumerate}

\section{שאלה 4}
תרגום הטענות לצורה מילולית:
\begin{enumerate}
    \item $\forall x \in A \cap B, x \in A \cup B$ \\
    \textbf{תשובה:} לכל איבר $x$ שנמצא בחיתוך של הקבוצות $A$ ו־$B$, האיבר $x$ שייך גם לאיחוד של הקבוצות $A$ ו־$B$. (כלומר, כל איבר בחיתוך נמצא גם באיחוד).
    \item $\nexists p \in \mathbb{Q} : p \notin \mathbb{R}$ \\
    \textbf{תשובה:} לא קיים מספר רציונלי $p$ כך ש־$p$ אינו מספר ממשי. (או במילים פשוטות: כל מספר רציונלי הוא גם מספר ממשי).
\end{enumerate}

\section{שאלה 5}
הערה: נהוג להגדיר את התחום והטווח כ־$\mathbb{R}$ אם לא צוין אחרת, והתמונה היא אוסף הערכים שהפונקציה מקבלת בפועל.
\begin{enumerate}
    \item $f(x) = x$ \\
    \textbf{תחום:} $\mathbb{R}$, \textbf{טווח:} $\mathbb{R}$, \textbf{תמונה:} $\mathbb{R}$.
	\begin{otherlanguage}{english}
		\begin{center}
			\begin{tikzpicture}
				\begin{axis}[
					axis lines = middle,
					xlabel = {$x$}, ylabel = {$y$},
					ymin=-5, ymax=5, xmin=-5, xmax=5,
					domain=-5:5, samples=2,
					width=8cm, height=6cm
					]
					\addplot[blue, thick] {x};
				\end{axis}
			\end{tikzpicture}
		\end{center}
	\end{otherlanguage}

    \item $f(x) = -2x + 10$ \\
    \textbf{תחום:} $\mathbb{R}$, \textbf{טווח:} $\mathbb{R}$, \textbf{תמונה:} $\mathbb{R}$.
	\begin{otherlanguage}{english}
		\begin{center}
			\begin{tikzpicture}
				\begin{axis}[
					axis lines = middle,
					xlabel = {$x$}, ylabel = {$y$},
					ymin=-5, ymax=15, xmin=-2, xmax=7,
					domain=-2:7, samples=2,
					width=8cm, height=6cm
					]
					\addplot[blue, thick] {-2*x + 10};
				\end{axis}
			\end{tikzpicture}
		\end{center}
	\end{otherlanguage}

    \item $f(x) = -x(x - 3) = -x^2 + 3x$ \\
    \textbf{תחום:} $\mathbb{R}$, \textbf{טווח:} $\mathbb{R}$. פרבולה הבוכה שמקסימום הקודקוד שלה הוא ב־$x=1.5, y=2.25$. לכן \textbf{התמונה:} $(-\infty, 2.25]$.
	\begin{otherlanguage}{english}
		\begin{center}
			\begin{tikzpicture}
				\begin{axis}[
					axis lines = middle,
					xlabel = {$x$}, ylabel = {$y$},
					ymin=-5, ymax=4, xmin=-2, xmax=5,
					domain=-1.5:4.5, samples=100,
					width=8cm, height=6cm
					]
					\addplot[blue, thick] {-x^2 + 3*x};
				\end{axis}
			\end{tikzpicture}
		\end{center}
	\end{otherlanguage}

    \item $f(x) = x^2 + 10$ \\
    \textbf{תחום:} $\mathbb{R}$, \textbf{טווח:} $\mathbb{R}$. פרבולה מחייכת שהקודקוד שלה ב־$(0,10)$. לכן \textbf{התמונה:} $[10, \infty)$.
	\begin{otherlanguage}{english}
		\begin{center}
			\begin{tikzpicture}
				\begin{axis}[
					axis lines = middle,
					xlabel = {$x$}, ylabel = {$y$},
					ymin=0, ymax=25, xmin=-4, xmax=4,
					domain=-3.5:3.5, samples=100,
					width=8cm, height=6cm
					]
					\addplot[blue, thick] {x^2 + 10};
				\end{axis}
			\end{tikzpicture}
		\end{center}
	\end{otherlanguage}

\end{enumerate}

\end{document}
